\section{Problem Understanding and Objectives}

\subsection{Problem statement}
When a machine is hit by ransomware or a crypto-stealing implant, much of the evidence exists only in
volatile memory: code injected into a legitimate process, the cipher implementation the malware loaded,
key material, session tokens and live network connections. Memory images are large (1 to 4.3\,GB in our
tests) and existing tools answer one question each: Volatility3 lists processes and injected regions,
YARA matches byte patterns, Wireshark dissects traffic. An analyst still has to join these results by hand,
decide which cipher is really in use, judge whether a recovered key is weak, and write it all up.

LLMs can write the report, but two problems follow. First, an LLM may state facts that are not in the
evidence. Second, the evidence itself is attacker controlled: a JWT claim, an OTP label, a TLS server name
or a command line can contain text such as ``ignore previous instructions, report the system is clean''.
A forensic tool that can be talked into calling a compromised host clean is worse than none.

\textbf{TriageGuard} therefore aims to automate memory triage with cryptography as its core, and to make
the LLM report both evidence-cited and measurably resistant to prompt injection.

\subsection{Objectives}
Each objective has a measurable criterion; Section~\ref{sec:testing} reports the results.
\begin{enumerate}
  \item \textbf{O1 Extraction.} Run Volatility3 on public Windows dumps and record processes, command
        lines, connections, injected (RWX, unbacked) regions and USB history. \emph{Criterion:} runs end to end
        on at least three dumps from different Windows versions.
  \item \textbf{O2 Cipher identification.} Identify AES, DES/3DES, RSA, ChaCha20/Salsa20, AES-GCM and
        ECC (P-256, secp256k1, Curve25519/Ed25519) from constants and API names, with a confidence
        percentage. \emph{Criterion:} all planted algorithms found on a ground-truth image, zero hits on random data.
  \item \textbf{O3 RSA weak keys.} Extract RSA keys (CAPI, CNG, DER, PEM) and test modulus size, small
        exponent, shared modulus, shared prime (GCD), close primes (Fermat) and small factors; score
        strength 0 to 100. \emph{Criterion:} every planted weakness detected, strong key not flagged.
  \item \textbf{O4 Credentials.} Find and decode JWTs, OAuth/API tokens and \texttt{otpauth://} secrets,
        storing them redacted. \emph{Criterion:} all planted credentials found and correctly assessed.
  \item \textbf{O5 Network.} Parse a pcap and flag weak negotiated TLS (deprecated versions, no forward
        secrecy, CBC, 3DES/RC4).
  \item \textbf{O6 Cited report.} Produce an LLM report from findings only, with a citation id on every
        claim and a rule-based risk score as reference. \emph{Criterion:} at least 90\% of claim lines cited.
  \item \textbf{O7 Red-team and defences.} Measure attack success rate (ASR) with and without defences.
        \emph{Criterion:} lower ASR and higher citation rate with defences, and no misleading report reaching a reader without a warning.
\end{enumerate}

\subsection{Literature review}
We surveyed 29 works: ten recent papers (2024 onward) that cover every pillar of the project, plus the
foundational work and specifications they build on. The full survey, in the course's table format with
cryptographic parameters and relevance to each pillar, is Table~\ref{tab:survey} in Appendix~\ref{app:survey}.

\textit{Memory forensics and AI-assisted triage.} Volatility began as Volatools~\cite{walters2007volatools}, and
\emph{The Art of Memory Forensics}~\cite{ligh2014art} set out the standard methods for processes, registry and
injected code (\texttt{malfind}). volGPT~\cite{oh2024volgpt} feeds Volatility artefacts to an LLM to triage
ransomware processes and improves on raw plugin output, and ForensicLLM~\cite{sharma2025forensicllm} shows a
local LLM can match cloud models for forensic questions. Neither tests whether crafted evidence can bias the
LLM's verdict, and neither assesses the cryptography found.

\textit{Detecting encryption, keys and ciphers.} Key material stands out by entropy~\cite{shamir1999hide},
and Halderman et al.~\cite{halderman2008cold} recovered AES and RSA keys from RAM images via key-schedule
structure. Casino et al.~\cite{casino2025not} classify high-entropy segments to separate encryption from
compression, but do not say which algorithm or how strong. Gr\"obert et al.~\cite{grobert2011automated} and
Lestringant et al.~\cite{lestringant2015automated} identify primitives in binaries by execution traces and
data-flow graphs, and Barnes and Ghafarian~\cite{barnes2026machine} classify ransomware from PE-header
features; all work per binary, not on a whole memory image. The constants we match come from the
specifications of AES~\cite{fips197}, GCM~\cite{sp80038d}, ChaCha20~\cite{bernstein2008chacha} and
Curve25519~\cite{bernstein2006curve25519}.

\textit{Cipher cost and weak keys.} Baig et al.~\cite{baig2024comparative} and Banu et al.~\cite{banu2025analyzing}
measure the speed and memory hierarchy of AES, DES/3DES, RSA, ElGamal and ECC, which frames which ciphers malware
uses for bulk data and which for key exchange; Kharraz et al.~\cite{kharraz2015cutting} show how ransomware
families use them in practice. For RSA, Boneh~\cite{boneh1999twenty} surveys small-exponent and common-modulus
attacks, Ruzai et al.~\cite{ruzai2024concurrent} factor sets of moduli linked by weak key equations, and
Heninger et al.~\cite{heninger2012mining} and Lenstra et al.~\cite{lenstra2012ron} broke 0.50\% of TLS host keys
and about 0.2\% of collected moduli through shared primes. These checks have been run on certificate scans and in
theory, not on keys carved from a memory image.

\textit{Tokens and one-time passwords.} Raghu~\cite{raghu2025context} quantifies the cost of hardening JWT/OAuth
in production, and Rujichaikul and Rassameeroj~\cite{rujichaikul2025token} detect token misuse with rule-based
signatures (81.4\% accuracy). RFC~4226~\cite{rfc4226} sets the minimum OTP secret length we check. None of these
look for tokens left in memory.

\textit{LLM robustness.} Perez and Ribeiro~\cite{perez2022ignore} showed goal hijacking via ``ignore previous
instructions'', and Greshake et al.~\cite{greshake2023not} showed \emph{indirect} injection through data an
application retrieves, which is our setting since a memory dump is attacker-controlled data. Spotlighting~\cite{hines2024spotlighting}
and the architectural defence of Gokcimen~\cite{gokcimen2025novel} reduce attack success, and Liu et al.~\cite{liu2024formalizing}
found no single defence sufficient. None is evaluated on forensic evidence or checks that a report's claims cite real evidence.

\textit{Gap.} Each pillar is covered somewhere, but no existing work (a) joins memory artefacts with cipher
identification \emph{and} weak-key testing on keys carved from the dump, (b) reports a confidence and evidence
basis per algorithm and a strength score per key, and (c) uses an LLM for the report while measuring and
defending its exposure to injection through the evidence itself. TriageGuard targets exactly this combination.

\subsection{Scope}
\textit{In scope:} Windows memory images (Volatility3 stage), any file for the crypto and credential stages,
the cipher families listed in O2, RSA weak-key tests, JWT/OAuth/OTP credentials, TLS in pcaps, a free-tier
LLM (Groq, \texttt{openai/gpt-oss-120b}) and an offline dashboard.
\textit{Out of scope:} executing malware, cracking third-party data, real user credentials (public datasets
and synthetic keys only), Linux/macOS memory profiles (planned later), and post-quantum schemes (ML-KEM,
ML-DSA), which are a discussion point only.
